% Appendix: exhaustive description of the MAST signals and their processing.

\section{Signals, processing, and data limitations}
\label{app:data}

This appendix documents every measured signal the model consumes, the exact
processing applied to it, and the limitations that processing does not
remove. It is deliberately more detailed than is customary, because in this
study the data-preparation choices --- not the model class --- turned out to
control the scientific conclusion.

\subsection{Provenance and access}

All data come from the UKAEA MAST open archive, Level-2 product, read as
Zarr groups from public object storage and converted to NetCDF per
discharge. Four groups are used: \texttt{summary}, \texttt{thomson\_scattering},
\texttt{equilibrium}, and the D$_\alpha$ extract of
\texttt{spectrometer\_visible}. Discharge metadata and free-text session-log
annotations come from the accompanying catalogue API and are used only for
the label audit (\S\ref{app:data-audit}), never as model input.

\subsection{The signals}
\label{app:data-signals}

\paragraph{Electron temperature and density profiles (Thomson scattering).}
The dominant measurement, and the sparsest. Each discharge provides
$T_{\mathrm{e}}$ and $n_{\mathrm{e}}$ on a $121$-point major-radius chord
sampled at $\SI{5}{ms}$ (typically $80$--$140$ time slices, median $83$
after screening). Values arrive in eV and m$^{-3}$; units are read from the
file attributes and verified, never assumed. Two properties dominate
everything downstream. First, the chord is horizontal at $Z = 0$ while the
magnetic axis of these discharges sits at $|Z_{\mathrm{axis}}|$ up to
$\SI{0.2}{m}$, so the chord never samples the true core: its closest
approach corresponds to $\rho \approx 0.3$. Second, and more restrictive,
the great majority of channels fail quality screening --- after processing,
the median discharge retains $10$ usable radial columns out of $65$ grid
points (range $4$--$16$), all at $\rho \gtrsim 0.7$. The model is therefore
identified on a thin outer annulus, and we make no claim about core
transport anywhere in this paper.

\paragraph{Actuator and global signals (summary).} Sampled at
$\SI{1}{ms}$ over roughly $[-0.1, 0.7]\,\mathrm{s}$ ($\sim$$790$ samples):
injected neutral-beam power $P_{\mathrm{NBI}}$ (W), plasma current $I_p$
(A, signed; the magnitude is used), line-averaged density (m$^{-3}$),
Ohmic power $P_{\Omega}$ (W), and radiated power $P_{\mathrm{rad}}$ (W).
$P_{\Omega}$ carries large transient spikes at breakdown and termination and
is clipped to $[0, 5]\,\mathrm{MW}$ before use; the raw trace is retained
separately as a diagnostic.

\paragraph{D$_\alpha$ emission.} Three photodiode channels sampled at
$\SI{20}{\micro s}$ --- $\sim$$35\,000$ samples per discharge, two orders of
magnitude faster than any other signal here. The channel sum is used. This
signal is \emph{never} a model input: it is used offline to construct weak
labels and as the target of the observation head. Its fast cadence matters
because the edge-localised-mode bursts that contaminate it are only a few
hundred microseconds wide and would be aliased away by any resampling to the
profile grid before filtering.

\paragraph{Equilibrium.} Poloidal flux map, magnetic-axis position, LCFS
contour, elongation, triangularity, stored energy, and $q_{95}$, at
$\SI{5}{ms}$. Their use is documented in full in
Appendix~\ref{app:equilibrium}.

\subsection{Processing chain}
\label{app:data-processing}

Three time bases coexist: D$_\alpha$ at $\SI{20}{\micro s}$, actuators at
$\SI{1}{ms}$, profiles at $\SI{5}{ms}$, and the equilibrium at $\SI{5}{ms}$.
Everything is resolved onto the profile grid, but only after any filtering
that must happen at native resolution.

\begin{enumerate}
\item \textbf{Radial mapping.} Each Thomson channel's major radius is mapped
to the flux label $\rho$ through the equilibrium
(Appendix~\ref{app:eq-rho}), then profiles are interpolated onto a uniform
$65$-point $\rho$ grid. Interpolation is constant-extrapolated outside the
measured range rather than zero-padded, because injecting
$T_{\mathrm{e}} = 0$ at an unmeasured radius creates a spurious gradient
that the diffusion operator would faithfully act on.
\item \textbf{Per-time-slice screening.} A slice is rejected unless it has a
minimum number of finite channels and a minimum radial span. Impulsive
outliers are removed by a rolling median-absolute-deviation test.
\item \textbf{Per-channel screening.} A radial channel is rejected unless
its median level, its bounded relative jump statistics, its relative span,
and its time coverage all pass fixed thresholds. A discharge is written only
if at least four stable channels survive.
\item \textbf{Representative density.} A single scalar density per time is
formed as a trimmed mean over trusted channels and broadcast over $\rho$,
rather than using the noisy density profile pointwise.
\item \textbf{Boundary trace.} The outermost surviving measured temperature
at each time becomes the Dirichlet edge condition, gap-filled in time.
\item \textbf{Actuator resampling.} Actuator traces are interpolated onto
the profile grid with endpoint carry. All model inputs are then divided by
\emph{fixed physical constants} (MW, MA, $10^{19}\,\mathrm{m}^{-3}$) --- never
by per-discharge statistics, so that the same actuator setting produces the
same model input in every discharge and at deployment.
\end{enumerate}

\subsection{Weak label construction}
\label{app:data-labels}

Labels come from D$_\alpha$ and are deliberately conservative.

\begin{enumerate}
\item \textbf{Lower envelope.} A rolling $15$th percentile over
$\SI{15}{ms}$ windows. Edge-localised modes are short positive spikes, so a
low quantile tracks the inter-burst baseline whose step-down is the actual
transition signature. A mean or median filter smears the bursts into the
baseline and destroys the step.
\item \textbf{Plasma gate.} Samples are considered only where $|I_p|$
exceeds $60\%$ of its 95th percentile and density exceeds $20\%$ of its
peak. This excludes the current ramp, where D$_\alpha$ is low simply because
recycling has not yet built up --- a low baseline there reflects the absence
of plasma, not improved confinement.
\item \textbf{Threshold.} A variance-maximising (Otsu) split of the gated
envelope distribution, rejected if the threshold sits outside the central
range, if either class is nearly empty, or if the class means differ by less
than $0.20$ of the normalised range.
\item \textbf{Acceptance rules.} A candidate H segment must persist for at
least $\SI{20}{ms}$; must be preceded by an L reference inside the gate;
and its entry must show a baseline step-down of at least $0.10$ of the
normalised range across $\SI{25}{ms}$ contrast windows. The last rule was
added after slow density-ramp discharges produced false positives by
drifting across the threshold.
\item \textbf{Exclusions.} Samples within a short window of every switch are
labelled ``transition'' and excluded from scoring, as are all out-of-gate
samples. Multiple H segments per discharge are permitted, so
back-transitions are labelled rather than suppressed.
\end{enumerate}

\subsection{Label audit and cohorts}
\label{app:data-audit}

Labels are compared against the free-text session-log annotations retrieved
for every discharge, giving three cohorts: \emph{confirmed transitions}
(extractor and log agree), \emph{confirmed negatives} (both indicate L
throughout), and \emph{ambiguous} (they disagree, or the log is silent).
On the present corpus these number \nTrans{}, \nNoTrans{}, and \nAmbig{}
respectively. The large ambiguous fraction is dominated by discharges whose
session-log entry is empty, not by genuine disagreement.

This has a direct consequence we state rather than work around: with only
\nNoTrans{} confirmed negatives, the false-positive rate of the model can be
estimated only weakly, and we report it as such. Ambiguous discharges are
never used to support a claim in either direction.

\subsection{Limitations of the data, stated}
\label{app:data-limits}

\begin{enumerate}
\item \textbf{Radial coverage.} Roughly $10$ of $65$ radial grid points
carry usable data, all outboard of $\rho \approx 0.7$. Core transport is
unidentified; the pedestal-region conclusion is the only one the data can
support.
\item \textbf{Temporal resolution mismatch.} The transition develops on a
$\sim$$\SI{20}{ms}$ scale while profiles are sampled every $\SI{5}{ms}$, so
only three or four profile samples resolve it. Event timing better than
$\sim$$\SI{10}{ms}$ is not measurable from the profile channel alone.
\item \textbf{Label--observation coupling.} The labels and the observation
head both derive from D$_\alpha$. The latent dynamics never see it, which
removes dynamical leakage, but the \emph{evaluation} remains statistically
tied to a single diagnostic. Independent confirmation would require an
electron-cyclotron or reflectometry-based transition marker, which the
Level-2 product does not provide.
\item \textbf{Midplane assumption.} The Thomson chord's vertical position is
absent from the files and is assumed to be $Z = 0$.
\item \textbf{Absorbed versus injected power.} Only injected beam power is
available; shine-through, orbit losses, and charge-exchange losses are not
resolved, so the loss-power proxy of \eqref{eq:ploss} systematically
overestimates the power actually crossing the separatrix. Thresholds derived
here are conditional, corpus-specific quantities and are not directly
comparable to published loss-power thresholds.
\item \textbf{Single device, narrow operating space.} All discharges are
MAST, in a limited current and density range. Nothing here tests transfer to
another device or to a substantially different scenario.
\item \textbf{Supervision-target normalisation.} The D$_\alpha$ target is
range-normalised per discharge. This affects training only --- no model input
is normalised this way and deployment needs no such statistic --- but it
means the observation head is fitted against a per-discharge-scaled target.
\end{enumerate}
